\section{Background and Related Work}
\label{sec:background}

\subsection{Dealer Gamma Exposure and 0DTE Options}
\label{sec:background:gex}

Market makers supply immediacy and carry the inventory that results \cite{grossman1988liquidity}. When they hedge option books toward delta-neutral, their hedge trades depend on the gamma of the book, and in models these trades feed back to the volatility of the underlying \cite{frey1997feedback}. In foreign-exchange options, aggregated market makers carry negative gamma that feeds back to the spot market \cite{anderegg2022impact}; in equities, option-related hedging has measurable price effects on the underlying \cite{ni2005stock,baltussen2021hedging}, and end-user demand moves option prices \cite{garleanu2009demand}. Practitioners summarise the dealers' position as a gamma exposure (GEX) figure \cite{spotgamma2021,fishman2023gamma}. The evidence that this position moves volatility is mixed: for S\&P~500 options, a recent study finds that same-day expiries dampen rather than amplify market volatility, through the hedging needs of market makers \cite{adams2025zerodte}. A GEX figure is also a lossy scalar: it depends on modelling choices (which gamma, which contracts, which expiries), so two daily GEX series with the same name need not agree (Section~\ref{sec:results:inputs} compares four). It rests on an assumption about positions as well: open interest records how many contracts are open but not which side dealers hold, so the convention used here (Equation~(\ref{eq:gex})) assigns dealers a long position in calls and a short position in puts.

The growth of zero-days-to-expiration (0DTE) options since 2022 concentrates gamma at very short maturities \cite{cboe2024zero,cboe2025spx0dte,dim2023odtes,adams2025zerodte}. Whether 0DTE contracts are included is one such modelling choice, and it can change the sign of the daily figure on some days. Regime-switching models detect regimes from observable outcomes such as returns \cite{hamilton1989new,ang2002regime,nystrup2020regime}, and model-free methods cluster windows of returns by the distance between their empirical distributions \cite{horvath2024clustering}. The rule and the LLM in this paper work on dealer positioning instead, and the regime they label is defined by the three thresholds of the rule; whether these windows coincide with regimes found from returns is not tested here.

\subsection{Validating LLM Reasoning via Obfuscation}
\label{sec:background:validation}

LLMs are applied to financial tasks such as return prediction from news \cite{lopez2023can} and multi-agent trading \cite{xiao2024tradingagents}, and chain-of-thought prompting makes their reasoning steps explicit \cite{wei2022chain,kojima2022large}. Three concerns motivate behavioral tests of the CheckList type \cite{ribeiro2020beyond}: accuracy depends on how probable the task, its input and its output are in the training text \cite{mccoy2023embers}; benchmark data can leak into training \cite{dong2024generalization}; and economic forecasts can reflect memorized history \cite{lopezlira2025memorization}. Temporal obfuscation removes dates, tickers and event context so that only numerical structure remains; it is one such test \cite{regan2025obfuscation}. In finance, the closest tests anonymize company names in news text: anonymized headlines can predict returns better in sample, which points to distraction rather than look-ahead bias \cite{glasserman2024lookahead}, and anonymizing also removes information that predicts outcomes \cite{wu2025anonymization}. Masking has a known limit: models can reconstruct masked entities and dates from little context \cite{lopezlira2025memorization}. The input here holds no text, but its scale still carries a clue to the period (Section~\ref{sec:discussion:threats}).

Obfuscation has a limited reach. It tests whether the answer depends on memorized context. When the prompt states the decision rule, as it does here, an answer that agrees with the rule shows that the model applies the stated criteria; it does not show how the model would answer without them. LLMs still struggle with open-ended reasoning about time series \cite{merrill2024timeseries}; the task here is narrower, the evaluation of three written criteria on 30 rounded values. Obfuscation also does not test whether the numbers shown to the model mean what the prompt says they mean. A model can pass an obfuscation test and still report faithfully on an input whose definition differs from what the prompt states. That second question concerns the inputs, not the model, and is the subject of the next subsection.

\subsection{Construct Validity of Financial Inputs}
\label{sec:background:construct}

An LLM validation inherits the construct validity of its inputs \cite{cronbach1955construct}. A daily GEX series is a derived quantity. Its definition (gamma model, contract filters, 0DTE treatment, put and call legs) and its provenance (vendor, load date, rebuild date) are part of the experiment, and they belong in the record of a study as much as the prompt and the model name do.

Three threats follow for the reproducibility of LLM studies. Data definitions drift between collections. Upstream data can be rebuilt silently, so that the same query returns a different series later. Models change and are deprecated: a hosted model can change its behavior between versions \cite{chen2024chatgpt}, settings meant to be deterministic still vary between runs \cite{atil2025nondeterminism}, and a retired model cannot be run again. In machine-learning-based science more broadly, leakage between inputs and labels is a common cause of results that do not reproduce \cite{kapoor2023leakage}. Section~\ref{sec:methods} describes how this paper records the input series and the model version for every arm.
